The System peripheral manages a potpourri of system-level blocks, including the MCU clocks, a CRC data verification module, analog bias levels within the global bias generator, and the bandgap voltage generator.

\subsection{System Clocks and Clock Management} \label{ss:clocks}
The MCU has three clock sources: HFXT, LFXT, and DCO0. HFXT and LFXT are external clock sources, which must be provided to the MCU via the \pin{ClkHFXT} and \pin{ClkLFXT} pins, respectively. DCO0 is an internal, programmable frequency clock source provided by a digitally controllable oscillator circuit.

HFXT must not exceed a maximum clock rate of about 100~\si{\mega\hertz}. LFXT is suggested to be exactly 32.768~\si{\kilo\hertz} to provide an accurate method of measuring time. The frequency of DCO0 is configured with the \register{DCO0FREQ} register, and can be up to 112~\si{\mega\hertz}.

Two clocks, the main clock (MCLK) and the submain clock (SMCLK), form the roots of the MCU clock tree, shown in Figure \ref{fig:mcu-clock-system}. MCLK and SMCLK can each select one of the source clocks as their base and divide the frequency further with \register{CLKDIVCR}. MCLK drives the CPU clock, while SMCLK drives most of the peripherals. MCLK cannot be turned off, but all other clocks and source clocks (including the CPU clock) can be turned off with the \register{SYSCLKCR} register. If a source clock is being used by MCLK, that source clock will be turned on (even if it is turned off in \register{SYSCLKCR}) unless the CPU clock is also turned off and an interrupt service routine is not executing. If a source clock is being used by SMCLK, it cannot be turned off with the \register{SYSCLKCR} register.

\begin{figure}[htpb]
\includesvg{figures/clocks-pingora2-2021-05}
\caption{MCU Clock System}
\label{fig:mcu-clock-system}
\end{figure}



\subsection{Power Saving Features} \label{ss:power-saving}
The MCU has several power-saving features.

The CPU is the most power-hungry part of the MCU, and its dynamic power consumption is proportional to the CPU clock frequency. Reducing the frequency of MCLK, the clock source for the CPU clock, is the main avenue to reduce the total dynamic power consumption. This can be done in one of two ways: either utilizing the MCLK clock divider, or directly reducing the frequency of the clock source. Reducing the clock source frequency is not always possible (such as when MCLK is sourced by an external clock on the \pin{ClkHFXT} pin), but it is always possible to switch the clock source to the internal DCO, which has a wide range of programmable frequencies to choose from.

Another method to save power is to turn off unused clocks, such as the internal DCO, when they are not being used. Remember, however, that the \pin{ClkHFXT} pin is required to be a valid clock signals while the MCU boots. Failure to provide the clock may cause the chip to boot improperly, or not at all. For this reason, the start HIGH pin(s) (\pin{SHx}) have been provided. On the PCB level, the start HIGH pins should connect to the enable pin of the clock generator chip for \pin{ClkHFXT} so that they will always be driven HIGH when the MCU is reset. After the chip has booted, the user can switch the source clock for MCLK and then safely disable \pin{ClkHFXT}.

The static power (or leakage power) can be reduced by switching off unused internal memories when they are not being used. Static power consumption, caused by the accumulated reverse bias currents of the many thousands of parasitic diodes in the digital circuits on the chip, is ever present as long as the chip is receiving power from the \pin{VDD}, \pin{VDDPST}, and \pin{AVDD} pins. The on-chip memory units, however, include a power gating feature that switches off the circuit delivering them power. The ROM and each 16~KiB SRAM block can be powered down by setting \bitfield{ROMOFF} and \bitfield{SRAMxOFF}, respectively. Every additional memory that is powered down will reduce the staic power consumption. Any read or write operations to a powered down will fail, producing undefined results. Furthermore, the entire contents of an SRAM block is lost when it is powered down, and is initially undefined when it is powered back on until it is filled again by the user. It is therefore important to keep track of which memories are or have been powered down to prevent undefined behavior.

When switching off SRAM blocks, it is recommended to power down the blocks at higher memory addresses first before powering down the lower blocks. The program occupies the lower memory addresses, which must not be corrupted by a negligent power down, but the stack, which is easily initialized in a custom location, starts at a high address and builds downward. To initialize the stack out of a powered down SRAM block and into a powered up block, add \bashinlinehl{-D RedefinedStackPointerInit=0x??} to the \bashinline{USER_DEFINES} variable in the project makefile, replacing \bashinline{0x??} with the desired memory address. The SRAM02 block should never be switched off because it contains the interrupt vector table.

The analog circuits on the chip can be powered down by disabling them in their control registers and disabling the global bias generator.



\subsection{CPU Sleep Mode} \label{ss:sleep}

Another way to conserve power is to only use the CPU when needed, and switch it off when not in use. As the CPU consumes the majority of the dynamic power on the chip, reducing the amount of time it is on will result in an overall power reduction. If properly configured, the peripherals can still operate even while the CPU is asleep, and can wake the CPU when certain events occur. Thus, the dynamic CPU power can be duty cycled such that the long term CPU power is equal to the percentage of the time that the CPU is on multiplied by the instantaneous power the CPU consumes when on: $P_\textrm{CPU} = P_\textrm{CPU, 100\%} \times t_\textrm{on}/(t_\textrm{on} + t_\textrm{off})$.

Sleep mode is activated with the custom \asminline{sleep} instruction, which reduces the dynamic power of the CPU to zero by switching off the CPU clock. Once it has entered sleep mode, the MCU can only awaken if an interrupt is triggered or if the MCU is reset. If an interrupt is triggered, the CPU clock is reactivated for the duration of the interrupt. Once all pending interrupts have been serviced, the MCU re-enters sleep mode. However, if one of the serviced ISRs executes the custom \asminline{wake} instruction, then the CPU will remain awake at the end of the master interrupt handler and jump back to the main part of the program where it left off when the CPU was put to sleep.

A code example to put the CPU to sleep is included in Listing \ref{lst:sleep}.

\begin{lstlisting}[style=customC, label=lst:sleep, caption={Example program to use CPU sleep mode}]
/** Includes **/
#include <MemoryMap.h>
#include <irq.h>

/** Interrupt Service Routine Declaration Macros **/
RVISR(IRQ_TIMER1_VECTOR, ISR_TIMER1)

/** Main Function **/
int main()
{
	// Set up TIMER1 here

	// Enable interrupts by unmasking them
	enable_all_interrupts();

	// Go to sleep
	cpu_sleep();

	// Insert code to execute after ISR wakes the CPU

	// Wait forever
	while (1) {}
}

/** Interrupt Service Routines **/
void ISR_TIMER1()
{
	// Do ISR processing here

	// Wake the CPU
	cpu_wake();
}
\end{lstlisting}



\subsection{CRC Module}
The CRC module enables the computation of a CRC16 checksum on an array of data. The checksum can be used to ensure the data in a transmission or reception is free of any corruption in the communication channel.

The specifications of the CRC module uses the CRC16 CDMA2000 standard, which has the following specifications:
\begin{itemize}
\item CRC size: 16 bits
\item Polynomial: 0xC857
\item Initial CRC value: 0xFFFF
\item No XOR operation on the output CRC value
\item No input data reflection
\item No output CRC value reflection
\end{itemize}

To use the CRC module on a byte array, first initialze the CRC value by writing 0xFFFF to \register{CRCSTATE}. Then, write each byte in the array to \register{CRCDATA}. Finally, read \register{CRCSTATE} to get the computed CRC16 value of the array.



\subsection{Global Analog Bias Generator}

The on-chip analog blocks receive bias voltages from the global bias generator circuit, which directly impacts the performance of all downstream analog blocks, such as the ADC, DACs, and opamps. Commercial chips usually go through rigorous calibration procedures to set the bias settings, which are then etched into the physical structure of the chip using laser engraving or one-time-programmable memories. This chip does not support one-time-programmable circuits, so the user needs to configure the global bias vector manually. Fortunately, this can be done using only the \register{BIASCR} register, which controls the bias levels for all analog circuits, and can disable the global bias generator and bandgap voltage generator as needed. Disabling the global bias generator also switches off all of the on-chip analog blocks, reducing their power consumtion such that they only consume leakage power.

To adjust the global bias current setting, modify \bitfield{BIASADJ}. Larger values of \bitfield{BIASADJ} yield larger currents, and smaller values yield smaller currents, albeit in a somewhat nonlinear fashion. The on-chip analog circuits are designed to have a nominal value of decimal 37 for \bitfield{BIASADJ}, but process variations, temperature changes, and power supply voltage deviations may require a different value of \bitfield{BIASADJ} to achieve the best results. The user should calibrate the on-chip analog circuits by sweeping \bitfield{BIASADJ} and measuring the analog performance at each test vector.

A schematic of the global bias generator is provided in Figure \ref{fig:bias-gen}.

\begin{figure}[htpb]
	\includesvg{figures/bias-generator-2020-05}
	\caption{Global Analog Bias Generator Schematic}
	\label{fig:bias-gen}
\end{figure}


\subsection{Global Bandgap Voltage Generator}

The chip also contains a global bandgap voltage generator, which is used as a stable, accurate voltage reference for the ADC and DACs. The bandgap voltage is nominally 1.196~V. The bandgap reference is resistant to process variations and temperature swings, with its output swinging from 1.191~V to 1.203~V across many doping concentration variations and temperatures from -25$^\circ$~C to 75$^\circ$~C. The power supply rejection ratio for the voltage is 34.5~dB. A schematic of the bandgap generator is shown in Figure \ref{fig:bandgap}.

\begin{figure}[htpb]
	\includesvg{figures/bandgap-generator-2021-05}
	\caption{Global Bandgap Voltage Generator Schematic}
	\label{fig:bandgap}
\end{figure}

% The explanation is from Razavi, Design of Analog CMOS Integrated Circuits, 1st edition
The bandgap reference is designed to output a DC voltage independent of the temperature. Unfortunately, most circuit components exhibit a nonzero temperature coefficient (TC), meaning that their voltages and currents are dependent on the temperature. The base-emitter junction voltage $V_{BE}$ of the bipolar transistors $Q_1$, $Q_2$, and $Q_3$ each have a negative TC. That is, as temperature increases, $V_{BE}$ decreases ($\partial V_{BE} / \partial T < 0$). The TC for base-emitter junctions is unfortunately nonlinear with respect to $T$. However, the voltage \textit{difference} between the base-emitter junctions of two bipolar transistors $\Delta V_{BE} = V_{BE1} - V_{BE2} = V_T \ln n$ is linear with respect to temperature, where $V_T = kT/q$ and $n = J_1/J_2$, the ratio of the current densities. The resulting temperature coefficient is therefore $\partial \Delta V_{BE} / \partial T = (k/q) \ln n$, which is a positive TC and not is dependent on the temperature.

The MOSFETs $M_{1-4}$ enforce identical currents on $Q_1$ and $Q_2$, despite their difference in emitter areas. $M_5$ mirrors that same current into $Q_3$. $Q_1$ and $Q_3$ are identically sized, which makes their $V_{BE}$ voltages and current densities equivalent. However, $Q_2$ has an emitter area $n$ times larger than the other two. As all three emitter currents are identical, the current density of $Q_2$ is thus $n$ times smaller than the others, giving rise to the equation $\Delta V_{BE} = V_T \ln n$, which has a positive TC. The MOSFETs also from an amplifier in feedback that ensures $V_{S2} = V_{S4}$, making the voltage across $R_1$ equal to $\Delta V_{BE}$. 

Thus, the output voltage is $V_{bg} = V_{BE3} + IR_2$. Note, however, that since the currents through $R_1$ and $R_2$ are equivalent and the voltage across $R_1$ is $\Delta V_{BE}$, the voltage across $R_2$ can be rewritten as $(R_2 / R_1) \Delta V_{BE}$. Hence, the output voltage of the bandgap circuit is $V_{bg} = V_{BE3} + (R_2 / R_1) \Delta V_{BE}$, which is nominally about 1.196~V.

To force the TC of the output voltage $V_{bg}$ to zero, the negative TC ($\partial V_{BE3} / \partial T$) needs to cancel the positive TC ($\partial \Delta V_{BE} / \partial T$). As the magnitude of $\partial V_{BE3} / \partial T$ is much larger than $\partial \Delta V_{BE} / \partial T$, $R_2 / R_1$ is sized such that $\partial \Delta V_{BE} / \partial T$ is amplified to be equal to the negative of $\partial V_{BE3} / \partial T$, yielding a zero TC for the output. Finally, it is important to note that while $\partial^2 \Delta V_{BE} / \partial T^2$ is equal to zero, $\partial^2 V_{BE3} / \partial T^2$ is nonzero, which makes the output TC equal to zero at only one temperature. However, this is not a strong influence on the output voltage, so the TC is about equal to zero for a wide range of temperatures centered on room temperature.
